% Part: computability
% Chapter: recursive-functions
% Section: primes

\documentclass[../../../include/open-logic-section]{subfiles}

\begin{document}

\olfileid[hi]{cmp}{rec}{pri}
\olsection{अभाज्य संख्याएँ}

परिबद्ध परिमाणीकरण और परिबद्ध न्यूनीकरण हमें यह दिखाने के लिए पर्याप्त
साधन देते हैं कि स्वाभाविक फलन और संबंध आदिम पुनरावर्ती हैं। मसलन,
संबंध ``$x$, $y$ को विभाजित करता है'' पर विचार करें, जिसे $x \mid y$
लिखते हैं। संबंध $x \mid y$ तब मान्य है जब~$y$ का $x$ से भाग शेष के
बिना संभव हो, अर्थात् जब $y$,~$x$ का कोई पूर्णांक गुणज हो। (यदि यह
मान्य नहीं है, अर्थात्~$x$ से $y$ को भाग देने पर शेष $> 0$ है, तो हम
$x \nmid y$ लिखते हैं।) दूसरे शब्दों में, $x \mid y$ तभी और केवल तभी
जब किसी~$z$ के लिए $x \cdot z = y$ हो। स्पष्टतः ऐसा कोई $z$, यदि वह
अस्तित्व में है, तो $\leq y$ होना चाहिए। अतः $x \mid y$ तभी और केवल
तभी जब किसी $z \le y$ के लिए $x \cdot z = y$ हो। हम संबंध $x \mid y$
को $=$ और गुणा से परिबद्ध अस्तित्वात्मक परिमाणीकरण द्वारा इस प्रकार
परिभाषित कर सकते हैं:
\[
x \mid y \defiff \bexists{z \leq y}{(x \cdot z) = y}.
\]
इस प्रकार हमने दिखाया कि $x \mid y$ आदिम पुनरावर्ती है।

कोई प्राकृतिक संख्या~$x$ \emph{अभाज्य} है यदि वह न $0$ है, न $1$, और
केवल $1$ तथा स्वयं से विभाज्य है। दूसरे शब्दों में, अभाज्य संख्याएँ ऐसी
हैं कि जब भी $y \mid x$ हो, तब या तो $y = 1$ अथवा~$y=x$। यह जाँचने के
लिए कि~$x$ अभाज्य है या नहीं, हमें केवल $y \mid x$ की जाँच उन सभी
$y \le x$ के लिए करनी है, क्योंकि यदि $y > x$, तो स्वतः~$y \nmid x$। अतः संबंध
$\fn{Prime}(x)$, जो तभी और केवल तभी मान्य है जब $x$ अभाज्य हो, इस प्रकार
परिभाषित किया जा सकता है:
\[
\fn{Prime}(x) \defiff x \geq 2 \land \bforall{y \leq x}{(y \mid x \lif y
  = 1 \lor y = x)}
\]
और इसलिए आदिम पुनरावर्ती है।

अभाज्य संख्याएँ $2$, $3$, $5$, $7$, $11$ इत्यादि हैं। उस फलन $p(x)$
पर विचार करें जो इस अनुक्रम की $x$वीं अभाज्य संख्या लौटाता है, अर्थात्
$p(0) =
2$, $p(1) = 3$, $p(2) = 5$ इत्यादि। (सुविधा के लिए हम प्रायः
$p(x)$ को $p_x$ लिखेंगे: $p_0=2$, $p_1=3$ इत्यादि।)

यदि हमारे पास ऐसा फलन $\fn{nextPrime(x)}$ हो जो~$x$ से बड़ी पहली
अभाज्य संख्या लौटाता है, तो~$p$ को आदिम पुनरावर्तन का उपयोग करके सरलता
से परिभाषित किया जा सकता है:
\begin{align*}
  p(0) & = 2\\
  p(x+1) & = \fn{nextPrime}(p(x))
\end{align*}
चूँकि $\fn{nextPrime}(x)$ वह लघुतम $y$ है जिसके लिए $y > x$ और~$y$
अभाज्य है, उसे अपरिबद्ध खोज द्वारा सरलता से संगणित किया जा सकता है। किंतु
यूक्लिड के एक परिणाम के कारण उसे परिबद्ध न्यूनीकरण द्वारा भी परिभाषित
किया जा सकता है: $x$ और $\fact{x}+1$ के बीच सदा कोई अभाज्य संख्या होती है।
\[
  \fn{nextPrime(x)} =
  \bmin{y \leq \fact{x}+1}{(y > x \land \fn{Prime}(y))}.
\]
इससे पता चलता है कि $\fn{nextPrime}(x)$ और इसलिए $p(x)$ (केवल संगणनीय
ही नहीं, बल्कि) आदिम पुनरावर्ती हैं।

(यदि आप जानना चाहें, तो यूक्लिड के प्रमेय का एक संक्षिप्त प्रमाण यह है।
मान लें कि $p_n$, $\le x$ वाली सबसे बड़ी अभाज्य संख्या है और गुणनफल
$p =
p_0\cdot p_1 \cdot \dots \cdot p_n$ पर विचार करें, जो सभी
अभाज्यों~$\le x$ का गुणनफल है। या तो $p+1$ अभाज्य है, या $x$ और~$p+1$ के बीच कोई
अभाज्य है। क्यों? मान लें कि $p+1$ अभाज्य नहीं है। तब कोई अभाज्य संख्या
$q \mid p+1$ ऐसी है कि $q < p+1$। अभाज्यों $\le x$ में से कोई भी
$p+1$ को विभाजित नहीं करता। (~$p$ की परिभाषा से प्रत्येक अभाज्य
$p_i \le x$,~$p$ को विभाजित करता है, अर्थात् शेष~$0$ रहता है। इसलिए
प्रत्येक अभाज्य $p_i \le x$,~$p+1$ को शेष~$1$ के साथ विभाजित करता है,
और अतः $p_i \nmid p+1$।) इसलिए $q$ ऐसा अभाज्य है जो $>
x$ और
$< p+1$ है। और $p \le \fact{x}$, इसलिए कोई अभाज्य $> x$ और $\le
\fact{x}+1$ है।)

\begin{prob}
परिबद्ध न्यूनीकरण का उपयोग करके पूर्णांक भाग $d(x, y)$ परिभाषित करें।
\end{prob}

\end{document}
